\section{Introduction}

The past two years have produced a curious convergence. Technical analyses document diminishing returns from scaling model size. Industry leaders publish essays calling for the ``pacing'' of frontier development. Regulators, largely absent from the AI conversation a short time ago, now face coordinated proposals for certification regimes, independent evaluators, and antitrust coordination. Each development is routinely read through one of two lenses: the \emph{technological} lens (AI has hit a wall) or the \emph{strategic} lens (incumbents are capturing regulation). Both capture something. Neither explains why the two developments are occurring together.

This essay argues that the convergence is best understood as a change in the economic gradient of deployment, not a slowdown in capability---and that the change matters because it alters what firms are competing for. The cost of achieving a fixed level of capability has fallen roughly thirteenfold per year since 2023~\cite{epochai2026}. The premium for frontier capability decays within months. Production traffic is migrating to commoditized open-weight models for high-volume workloads, while frontier labs retain high-value, reliability-sensitive tasks~\cite{vercel2026}. Value accrues elsewhere; the frontier's ceiling is not the issue. When differentiation becomes harder to monetize, control over the infrastructure layer---distribution, compute, and the governance arrangements that determine access---becomes the site of rent.

The essay makes three moves, and it is important to distinguish them because they carry different evidentiary weight. First, it establishes the empirical shift using production data---Vercel's AI Gateway Index and Epoch AI's inference-cost analysis---rather than benchmarks. This is the empirical finding. Second, it interprets the shift through Perez's framework, treating the installation/deployment distinction as a source of testable expectations rather than a retrospective label. This is a theoretical interpretation, not a demonstration. Third, it evaluates the institutional response---the Anthropic/METR arrangement and its legislative generalization---through an eight-dimension framework that distinguishes observed design features from mechanisms of exclusion and from demonstrated effects. This is the analytical contribution, and it is independent of the causal interpretation.

The methodological stance is explicit and load-bearing: \textbf{this essay does not adjudicate intent.} A regime designed with the best of intentions can produce exclusionary effects; a regime designed with cynical motives can produce genuine safety benefits. What matters is what the design \emph{does}---who is subject to it, who controls it, who bears its costs, who can contest its findings. The framework is therefore a multidimensional audit, falsifiable per dimension, with verdicts reported per-dimension. If Anthropic's draft bill, when introduced, includes government funding, shared infrastructure, proportionality provisions, and an open-weight regime, the exclusionary prediction is disconfirmed on those dimensions. The essay asks the reader to evaluate the design, not to take a position on the motives of the firms involved.

A note on what the essay does not claim. It does not claim that AI has reached fundamental limits. It does not claim that a Perez-style phase transition has definitively occurred. It does not claim that frontier laboratories are coordinating to suppress competitors. It claims that the economic gradient is shifting (empirical), that the shift is consistent with a transition from product to infrastructure competition (theoretical), that the shift may increase the incentive to shape governance (hypothesis, one of three non-exclusive explanations), and that the design features of the governance response have distributional consequences regardless of which explanation is correct (analytical contribution). These claims are separated throughout.

The essay proceeds as follows. Section~2 sets out the Perez framework and derives explicit expectations. Section~3 establishes the empirical shift, with the C$\rightarrow$B mechanism (falling inference cost drives deployment migration) as the anchor. Section~3.5 introduces the distinction between product logic and infrastructure logic, and defines the attractor heuristically. Section~4 examines the Anthropic/METR arrangement through eight questions, establishing what the evidence shows and what it does not, and presents three non-exclusive explanations for the institutional response. Section~5 presents the eight-dimension evaluative framework as a multidimensional audit, expands its relationship to existing AI-governance frameworks (the EU AI Act, the NIST AI RMF, the UK AI Safety Institute), and validates the framework against the EU AI Act's conformity-assessment regime, China's algorithm registration regime, and the UK AI Safety Institute. Section~6 tests the framework against four historical cases and states where each analogy fails. Section~7 concludes with what is established, what remains uncertain, and a research agenda.

The question is not whether AI will continue to develop. It will. The question is whether the institutional arrangements being constructed now will enable broad-based gains or entrench the position of those who built the infrastructure. The design choices are being made now, and incumbent firms are playing a significant role in making them. That is the fact the essay establishes. What follows from it is a matter of design, not destiny.
